\section{Compacidad}

% Capítulo 20: Compacidad en espacios métricos

\subsection{Número de Lebesgue}

\begin{definición}[Número de Lebesgue]
    Sea $(X,d)$ métrico y $\mathcal U$ un recubrimiento abierto de $X$. Un número $\delta>0$ es un \textbf{número de Lebesgue} de $\mathcal U$ si toda bola $B(x,\delta)$ queda contenida en algún elemento de $\mathcal U$.
\end{definición}

\begin{teorema}
    Si $(X,d)$ es compacto, todo recubrimiento abierto de $X$ admite número de Lebesgue.
\end{teorema}

\subsection{Totalmente acotado y completo}

\begin{definición}[Totalmente acotado]
    Un espacio métrico $(X,d)$ es \textbf{totalmente acotado} si para todo $\varepsilon>0$ existe un recubrimiento finito de $X$ por bolas de radio $\varepsilon$.
\end{definición}

\begin{proposición}
    Todo compacto métrico es totalmente acotado y completo.
\end{proposición}

\begin{teorema}[Caracterización métrica]
    En espacios métricos:
    $$X \text{ compacto } \Longleftrightarrow X \text{ completo y totalmente acotado.}$$
\end{teorema}

% Capítulo 20: Espacios compactos

\subsection{Recubrimientos}

\begin{definición}[Recubrimiento abierto y compacto]
    Un recubrimiento abierto de $X$ es una familia $\mathcal U\subset \tau$ con
    $$X=\bigcup_{U\in\mathcal U} U.$$
    El espacio $(X,\tau)$ es \textbf{compacto} si todo recubrimiento abierto admite subrecubrimiento finito.
\end{definición}

\begin{proposición}[Caracterización por cerrados]
    $X$ es compacto si y sólo si toda familia de cerrados con la propiedad de intersección finita tiene intersección no vacía.
\end{proposición}

\subsection{Propiedades principales}

\begin{proposición}
    Toda imagen continua de un compacto es compacta.
\end{proposición}

\begin{corolario}
    La compacidad es invariante topológica.
\end{corolario}

\begin{proposición}
    En espacios Hausdorff, los subconjuntos compactos son cerrados.
\end{proposición}

\subsection{Producto de compactos}

\begin{teorema}[Tychonoff finito]
    Si $X$ e $Y$ son compactos, entonces $X\times Y$ es compacto.
\end{teorema}
